% !TeX program = xelatex
% 2026 CUMCM 电子版论文：第一页必须是摘要专用页。
% 编译：xelatex paper_v1.tex -> bibtex paper_v1 -> xelatex x2
% 本文件由 main.tex + sections/*.tex 自动合并生成；修改请改 sections 源文件后用 _merge.py 重新生成
\documentclass[zihao=-4,a4paper,UTF8,fontset=none]{ctexart}

% template
\usepackage{cumcm2026}

% set English and Number font
\setmainfont{TimesNewRoman.ttf}[
  Path=fonts/,
  BoldFont=TimesNewRomanBold.ttf,
  ItalicFont=TimesNewRomanItalic.ttf,
  BoldItalicFont=TimesNewRomanBoldItalic.ttf
]

% set Chinese font
\setCJKmainfont[
  Path=fonts/,
  BoldFont=SimHei.ttf
]{SimSun.ttc}

\setCJKfamilyfont{zhsong}[
  Path=fonts/
]{SimSun.ttc}

\setCJKfamilyfont{zhhei}[
  Path=fonts/,
  AutoFakeBold=2.5
]{SimHei.ttf}

% 代码环境(listings \ttfamily)中的中文使用宋体
\setCJKmonofont[
  Path=fonts/
]{SimSun.ttc}

\providecommand{\songti}{}
\providecommand{\heiti}{}
\renewcommand{\songti}{\CJKfamily{zhsong}}
\renewcommand{\heiti}{\CJKfamily{zhhei}}

%--------------------------------------------------------------------
% 标题与关键词（团队最终确认）
\title{基于储能套利与滚动优化的微网电力调控策略研究}
\newcommand{\keywords}{微网调度；储能套利；线性规划；随机优化；动态规划}

\begin{document}

% 2026 电子提交版不得包含承诺书和编号专用页；第一页直接为摘要页。
\maketitle

\begin{cumcmabstract}
% 摘要为全队定稿部分：Q1 内容已按冻结结果写入，Q2--Q4 部分待对应问题求解完成后补充修改。
针对微网与外部电网的电力调控问题，本文围绕“购电量与储能充放电量的最优调度”这一核心任务，构建了从确定性日前优化到不确定性滚动调整的递进式优化模型体系。

针对问题一，在电价、小区负载均已知且储能日初日末电量相等的确定性条件下，建立了以全天购电费用最小为目标、包含电力平衡、储能状态转移、容量与功率限制及边界条件的线性规划模型；利用储能“低价充电、高价放电”的套利机理，通过 HiGHS 求解器求得全局最优解：全天购电量为 $59{,}482.70$ kWh，全天购电费为 $35{,}126.95$ 元。与不使用储能的基线方案（$48{,}052.05$ 元）相比，储能运行带来净节省 $12{,}925.10$ 元，节省比例达 $26.9\%$。逐时段约束验证表明模型解严格满足全部约束，能量守恒误差为 $0.0000$ kWh，验证了模型与求解的正确性。

针对问题二，在负载与光伏逐日波动的条件下，构建“日前随机优化 + 日内动态规划执行”的两阶段购电模型：日前层基于历史联合残差情景求解两阶段随机线性规划，锁定当天计划购电量；日内层以历史路径平均价值函数为引导，按因果规则实时调度储能充放电，缺口按 5 倍电价紧急购电补齐。对 2025.2.1--12.31 共 334 天的回测显示：年度总购电费用为 $14{,}120{,}324$ 元（计划购电 $91.1\%$、紧急购电 $8.9\%$），相比无储能方案节省 $2{,}286{,}996$ 元（$13.94\%$）；动态规划执行器相比固定计划基准改进 $14.31\%$（$p<0.0001$），跨日连续性误差小于 $10^{-6}$ kWh，储能效率 $81.00\%$ 与理论值吻合。

针对问题三，构建四阶段确定性滚动优化模型：在 0:00 制定日前计划，并在 6:00、12:00、18:00 根据新光伏预报判断是否重优化剩余时段的购电策略，调整部分按题目规则计费（减购退 $50\%$、增购 $1.5$ 倍）。334 天回测结果为：全年总费用 $16{,}489{,}545$ 元（计划购电 $90.8\%$、紧急购电 $7.4\%$、调整费用 $1.9\%$），$41.0\%$ 的日期触发策略调整，紧急购电占比 $2.23\%$；能量守恒审计（含弃光项）误差为 $0.000000$ kWh。分析表明，引入多时点预报的滚动调整策略是必要且经济的。

针对问题四，将电价输入替换为附件4 的逐日逐时段波动电价，其余决策结构与约束体系完全继承问题二、问题三。波动电价下重算问题二（Q4-2）的 334 天回测结果为：年度总购电费用 $14{,}777{,}634.02$ 元（计划购电 $91.2\%$、紧急购电 $8.8\%$），较固定电价情形增加 $4.7\%$，成本结构保持稳定；重算问题三（Q4-3）的年度总运营成本为 $18.07$ 百万元（计划 $97.8\%$、调整 $1.0\%$、紧急 $1.2\%$），较问题三增加 $9.6\%$，紧急购电日均占比 $0.39\%$。储能与紧急购电机制在波动价格环境下依然有效，全部约束零违反，求解成功率 $100\%$。

\cumcmkeywords{\keywords}
\end{cumcmabstract}

% 问题重述
\section{问题重述}

\subsection{问题背景}

微网是为促进光伏等分布式能源就地消纳、缓解并网压力而发展起来的小型发用电系统。非孤岛运行的微网可以与外部电网进行电力交易：当外网电价较低或分布式电源出力富余时，为储能设备充电；当外网电价较高时，释放储存的电能。微网调控的核心任务是，在小区负载、光伏出力与储能电量等信息的支撑下，给出各时段的购电量与储能充放电量，使微网供电始终满足小区负载需求，同时尽可能降低购电费用。

某小区微网由光伏发电系统、储能设备与小区负载构成。储能设备最大容量为 $12000$ kWh，最大充放电功率为 $5000$ kW，充放电效率为 $90\%$，其电量须保持在 $1200$--$10800$ kWh 之间；2025年1月1日 0:00 的初始电量为 $6000$ kWh。题目提供了小区一天内以 10 分钟为间隔的电价、负载与光伏功率数据（附件1），以及全年的实际运行与预报数据（附件2--4），要求制定不同条件下的购电策略。

\subsection{问题内容}

\textbf{问题一：} 在每天电价和小区负载相同、光伏发电功率预测已知的确定性条件下，要求微网提供的电能不低于小区负载，且储能设备在 0:00 与 24:00 的储电量相同，建立每天 0:00 制定当天计划购电策略的数学模型；基于附件1与附录1的数据，给出指定时间段的购电量、全天购电量与购电费（表1），以及储能设备分时段充放电量与边界储电量（表2），并将完整策略写入 result1.xlsx。

\textbf{问题二：} 在每天电价相同、负载与光伏功率逐日变化的条件下，若微网供电低于负载，需向外网紧急购电（电价为交易时刻电价的 5 倍），建立每天 0:00 的计划购电模型，并基于附件1、附件2给出全年策略及指定日期的结果。

\textbf{问题三：} 在每天 0:00、6:00、12:00 与 18:00 可获取未来 24 小时光伏功率预报的条件下，允许根据新预报调整购电策略（调减部分收取 $50\%$ 违约费用，调增部分按 $1.5$ 倍电价计费），建立含计划、紧急与调整费用的购电模型，并分析引入滚动调整的必要性。

\textbf{问题四：} 在电价实时波动的条件下，重建购电策略模型，并在波动电价下重新求解问题二与问题三。

四个问题层层递进：问题一建立确定性日前优化的基准模型；问题二引入实际运行中的供需偏差与紧急购电机制；问题三进一步引入多时点预报与策略调整；问题四将确定电价推广为波动电价。本文依次建模求解，并在最后检验模型的有效性与稳健性。

% 问题分析
\section{问题分析}

\subsection{数据理解}

附件1 提供了某一天以 10 分钟为间隔的电价、小区负载与光伏发电预测功率数据，共 144 个时段（其中 ``0:00+1'' 表示当天 24:00）。对附件1 进行统计分析，结果如表~\ref{tab:data_stat} 所示。

\begin{table}[H]
  \centering
  \caption{附件1 主要变量的描述性统计}
  \label{tab:data_stat}
  \begin{tabular}{ccccccc}
    \toprule
    变量 & 均值 & 最小值 & 最大值 & 峰谷比/波动 & 单位 \\
    \midrule
    电价 & 0.766 & 0.371 & 1.395 & 3.76 & 元/kWh \\
    小区负载 & 4626.0 & 3309.4 & 5959.0 & 1.80 & kW \\
    光伏发电预测功率 & 2311.8 & 0.0 & 7612.3 & -- & kW \\
    \bottomrule
  \end{tabular}
\end{table}

数据呈现三个明显特征。电价具有显著的分时结构：最低电价 $0.371$ 元/kWh 与最高电价 $1.395$ 元/kWh 相差约 3.76 倍，为储能“低价充电、高价放电”的套利运行提供了经济空间。小区负载的日间波动同样明显（$3309$--$5959$ kW），存在早晚两个用电高峰。光伏发电则集中在白天：全天 89 个时段出力为正（占 $61.8\%$），峰值功率 $7612.3$ kW 超过负载峰值，光伏富余时段需要储能吸收多余电量。此外，数据完整无缺失、数值均在合理范围，可直接用于建模。

\subsection{总体思路}

四个子问题共享同一个物理系统与调控目标，区别在于信息条件与不确定性逐问增强，本文据此构建递进式的优化模型体系：

\textbf{问题一}是确定性日前优化：电价、负载与光伏出力均精确已知，且要求储能日初日末电量相等、实现日循环运行。其本质是一个储能套利调度问题，目标函数与全部约束均为决策变量的线性函数，故建立线性规划模型，借助成熟的单纯形类求解器即可获得全局最优解。该模型同时作为后续问题的基准框架。

\textbf{问题二}在问题一的基础上引入负载与光伏的实际波动以及“供电不足时按 5 倍电价紧急购电”的惩罚机制，需在日前计划中考虑不确定性的代价。

\textbf{问题三}引入 0:00、6:00、12:00、18:00 四个预报时刻，将问题转化为含计划、调整与违约费用的滚动优化问题。

\textbf{问题四}将确定电价替换为实时波动电价，检验前述模型在更一般价格信号下的适应性。

四个问题之间的逻辑关系与数据流如图~\ref{fig:roadmap} 所示（技术路线图待 A/B 确认后由 DrawIO 生成并插入）。

\begin{figure}[H]
  \centering
  % TODO: 技术路线图由 4drawio 阶段生成后插入，此处暂以文字占位，不占用最终图号
  \fbox{\parbox{0.8\textwidth}{\centering 技术路线图（待插入）}}
  \caption{总体技术路线}
  \label{fig:roadmap}
\end{figure}

\subsection{问题一的分析}

问题一的决策内容是全天 144 个时段的购电量与储能充放电量。三个环节相互耦合：\textbf{电力平衡}要求任一时刻购电、光伏与储能放电之和覆盖负载与储能充电需求；\textbf{储能动态}刻画充放电对储电量的改变及其效率损耗；\textbf{经济目标}要求在满足供电的前提下最小化购电费用。储能的存在使各时段决策相互关联——当前时段充入的电量可以在后续高价时段释放，将购电需求从高价时段“转移”到低价时段，这正是储能套利的机理。因此问题一可以表述为一个多时段耦合的线性规划问题：以各时段购电量与充放电量为决策变量，以全天购电费用为目标，以电力平衡、储能状态转移、容量与功率限制、边界条件为约束，求解器给出的即为理论最优的日前调度策略。

依据题意，问题一还需满足若干结构性条件：储能电量运行区间为 $1200$--$10800$ kWh（附录1 硬性要求）；充放电效率 $90\%$ 按字面理解为充电、放电单侧效率各为 $90\%$，即往返效率 $81\%$；0:00 与 24:00 储电量均为初始值 $6000$ kWh；微网不允许向外网售电，购电量非负，光伏富余时允许弃光。

% 模型假设
\section{模型假设}

本文在建模过程中采用以下假设：

\noindent 1) \textbf{确定性信息假设（问题一）}：一天内各时段的电价、小区负载与光伏发电功率均精确已知，且 10 分钟内保持不变。附件1 以 10 分钟为间隔给出数据，该假设与数据粒度一致，是日前优化模型的通用前提。

\noindent 2) \textbf{效率口径假设}：储能充放电效率按题目表述取单侧效率 $\eta = 90\%$，即充电输入 $E_{\mathrm{charge}}$ 时仅 $\eta E_{\mathrm{charge}}$ 存入电池，电池放出 $E_{\mathrm{discharge}}$ 时仅 $\eta E_{\mathrm{discharge}}$ 供给负载，往返效率为 $\eta^2 = 81\%$。效率在问题一的时间尺度内视为恒定，不随充放电功率与荷电状态变化。

\noindent 3) \textbf{储能物理约束假设}：储能设备同一时段不同时充电与放电。由于购电价格为正面充放电存在损耗，同时充放电只会产生净损耗而不降低购电费用，最优解将自然规避该情形，故模型中无需显式加入互斥约束（后文给出严格论证）。

\noindent 4) \textbf{电量平衡与边界假设}：微网供电允许大于负载需求，即光伏出力超过“负载＋储能可吸收电量”时允许弃光，弃光不产生收益也不产生成本；储能设备在 0:00 与 24:00 的储电量均为 $6000$ kWh，实现日循环运行。

\noindent 5) \textbf{交易规则假设}：微网只能从外网购电，不向外网售电；购电费用等于购电量与对应时段电价的乘积，无其他附加费用；问题一暂不考虑储能设备折旧、电网功率约束与电能质量要求等次要因素。

\noindent 6) \textbf{数据有效性假设}：赛题提供的附件数据真实有效，单位统一，无缺失与异常值（经数据审查验证）。

% 符号说明
\section{符号说明}

本文用到的主要符号及其含义如表~\ref{tab:symbols} 所示。符号体系统一在全部子问题中沿用，仅问题三、问题四增加少量带上下标的变量（届时在对应章节补充说明）。

\begin{table}[H]
  \centering
  \caption{主要符号与定义}
  \label{tab:symbols}
  \begin{tabularx}{0.92\textwidth}{C{0.22\textwidth}Y C{0.14\textwidth}}
    \toprule
    符号 & 含义 & 单位 \\
    \midrule
    $t$ & 时段索引，$t=1,2,\ldots,144$，对应 10 分钟间隔 & -- \\
    $\Delta t$ & 时间间隔长度，$\Delta t=1/6$ & h \\
    $Price(t)$ & 时段 $t$ 的电价 & 元/kWh \\
    $Load(t)$ & 时段 $t$ 的小区负载功率 & kW \\
    $PV(t)$ & 时段 $t$ 的光伏发电预测功率 & kW \\
    $E_{buy}(t)$ & 时段 $t$ 从外网购买的电量 & kWh \\
    $E_{charge}(t)$ & 时段 $t$ 的名义充电能量（输入侧，入池 $\eta E_{charge}$） & kWh \\
    $E_{discharge}(t)$ & 时段 $t$ 的名义放电能量（电池侧，供负载 $\eta E_{discharge}$） & kWh \\
    $E_{storage}(t)$ & 时刻 $t$ 结束时储能设备的储电量 & kWh \\
    $E_{\min}$, $E_{\max}$ & 储能电量下限、上限（$1200$, $10800$） & kWh \\
    $P_{\max}$ & 储能最大充放电功率（$5000$） & kW \\
    $\eta$ & 储能单侧充放电效率（$0.9$，往返 $0.81$） & -- \\
    $Z$ & 全天购电费用（目标函数值） & 元 \\
    \bottomrule
  \end{tabularx}
\end{table}

\noindent 说明：时段内变化的量（购电量、充放电量）以下标 $t$ 对应时段编号；时刻定义的量（储电量）对应时间点编号，$E_{storage}(0)$ 与 $E_{storage}(144)$ 分别表示 0:00 与 24:00 的储电量。功率与电量的换算关系为电量 $=$ 功率 $\times \Delta t$。

% 问题一：确定性日前购电优化模型
\section{问题一：确定性条件下的日前购电优化模型}

\subsection{建模思路}

问题一在电价、小区负载与光伏出力全部精确已知的条件下，求解全天 144 个时段的购电量与储能充放电量，使购电费用最小。这是典型的微网确定性经济调度问题，其决策结构与光伏--电池混合系统的需求侧管理问题一致\cite{wu2015demand}：储能设备扮演能量“时间搬运”的角色，在低价时段吸收外网电量与光伏富余电量，在高价时段释放，从而将购电需求从高价时段转移到低价时段。

该问题的目标函数（购电费用为购电量与电价的乘积之和）与全部约束（电力平衡、储能状态转移、容量与功率限制）均为决策变量的线性函数，且决策变量为连续变量，因此可直接建立线性规划（LP）模型。微网能量管理领域的综述研究表明，对于确定性、单目标的经济调度问题，线性规划凭借其全局最优性和求解效率是公认的首选方法\cite{zia2018microgrids}，无需引入元启发式算法。问题一的最优解即为微网在确定性信息下可达到的理论最优日前策略，同时也是问题二至问题四建模的基准框架。

\subsection{数据预处理}

附件1 提供了一天 144 个时段的电价、小区负载与光伏发电预测功率，时间间隔 $\Delta t = 10\ \mathrm{min} = 1/6$ h，表中 ``0:00+1'' 即当天 24:00。预处理步骤如下：

\textbf{(1) 时间索引化。} 将 144 个时段依次编号 $t = 1,2,\ldots,144$，其中 $t=1$ 对应 0:00--0:10，$t=61$ 对应 10:00--10:10，$t=144$ 对应 23:50--24:00。

\textbf{(2) 数据完整性校验。} 检查确认附件1 共 144 行、字段齐全，电价、负载与光伏功率均无缺失、无穷或非法数值，电价严格为正，负载与光伏功率非负，时间序列自 0:00 起严格等间隔 10 分钟且顺序正确。

\textbf{(3) 单位统一。} 模型以电量为计量单位，将功率数据乘以时间间隔换算为电量：$Load(t)\cdot\Delta t$、$PV(t)\cdot\Delta t$ 分别表示时段 $t$ 的负载需求电量与光伏发电电量（kWh）。储能每时段最大充放电电量为 $P_{\max}\Delta t = 5000/6 = 833.33$ kWh。

\subsection{储能充放电的线性规划模型}

\subsubsection{决策变量}

对每个时段 $t\in\{1,2,\ldots,144\}$ 定义四类连续决策变量：

\begin{itemize}
  \item $E_{buy}(t)$：时段 $t$ 从外网购买的电量（kWh）；
  \item $E_{charge}(t)$：时段 $t$ 储能的名义充电能量（kWh），即充电输入侧能量，经效率 $\eta$ 折算后实际入池 $\eta E_{charge}(t)$；
  \item $E_{discharge}(t)$：时段 $t$ 储能的名义放电能量（kWh），即电池放出侧能量，经效率 $\eta$ 折算后实际供负载 $\eta E_{discharge}(t)$；
  \item $E_{storage}(t)$：时刻 $t$ 结束时储能的储电量（kWh），$t=1,\ldots,144$。
\end{itemize}

共 $144\times 4 = 576$ 个决策变量；$E_{storage}(0)=6000$ kWh 为已知参数。效率 $\eta=0.9$ 统一作用于能量流出方向：充电时输入 $E_{charge}$ 损失 $10\%$、入池 $\eta E_{charge}$；放电时电池放出 $E_{discharge}$ 损失 $10\%$、供负载 $\eta E_{discharge}$，往返效率为 $\eta^2=0.81$。

\subsubsection{目标函数}

以全天购电费用最小为目标：

\begin{equation}
  \min Z = \sum_{t=1}^{144} Price(t)\cdot E_{buy}(t)
  \label{eq:obj}
\end{equation}

光伏发电与储能放电均不产生购电成本，模型将优先利用这两类零成本能量。

\subsubsection{约束条件}

\textbf{(1) 电力平衡约束。} 任一时刻微网供给能力须不低于负载需求：

\begin{equation}
  E_{buy}(t) + PV(t)\Delta t + \eta E_{discharge}(t) \ge Load(t)\Delta t + E_{charge}(t),\quad t=1,\ldots,144
  \label{eq:balance}
\end{equation}

其中左侧为供给（购电、光伏发电、储能放电），右侧为需求（负载用电、储能充电输入）。约束取不等式，允许光伏出力超过“负载与储能吸收能力之和”时弃光；由于目标只惩罚购电量且无售电收益，最优解不会故意产生多余供给。

\textbf{(2) 储能状态转移方程。} 相邻时刻储电量满足：

\begin{equation}
  E_{storage}(t) = E_{storage}(t-1) + \eta E_{charge}(t) - E_{discharge}(t),\quad t=1,\ldots,144
  \label{eq:state}
\end{equation}

\textbf{(3) 容量约束。} 储电量须始终处于题目规定的安全运行区间：

\begin{equation}
  1200 = E_{\min} \le E_{storage}(t) \le E_{\max} = 10800,\quad t=1,\ldots,144
  \label{eq:capacity}
\end{equation}

\textbf{(4) 充放电功率约束。} 名义充放电能量受最大充放电功率限制：

\begin{equation}
  \begin{split}
    0 \le E_{charge}(t) \le P_{\max}\Delta t = 833.33,\quad t=1,\ldots,144,\\
    0 \le E_{discharge}(t) \le P_{\max}\Delta t = 833.33,\quad t=1,\ldots,144
  \end{split}
  \label{eq:power}
\end{equation}

\textbf{(5) 边界条件。} 储能日初、日末电量相同，实现日循环运行：

\begin{equation}
  E_{storage}(0) = 6000,\qquad E_{storage}(144) = 6000
  \label{eq:boundary}
\end{equation}

\textbf{(6) 非负约束。} 微网不向外网售电：

\begin{equation}
  E_{buy}(t) \ge 0,\quad t=1,\ldots,144
  \label{eq:nonneg}
\end{equation}

\subsubsection{充放电互斥性的理论论证}

模型式\eqref{eq:balance}--\eqref{eq:nonneg}为纯线性规划，未显式禁止同时充放电。以下证明最优解不会出现同时充放电，故无需引入 0-1 互斥变量。设某可行解在时段 $t$ 同时有 $E_{charge}(t)=c>0$、$E_{discharge}(t)=d>0$，令 $\delta=\min(c,d)>0$，构造修正解 $\tilde c=c-\delta$、$\tilde d=d-\delta$（其余变量不变）。修正后：电力平衡左侧减少 $\eta\delta$、右侧减少 $\delta$，由 $\eta<1$ 知约束\eqref{eq:balance}的松弛量增加，仍然成立；状态方程\eqref{eq:state}右端增加 $(1-\eta)\delta>0$，储电量增大，容量约束\eqref{eq:capacity}不会因此被破坏；目标函数值不变。故任一含同时充放电的可行解均可经上述修正消去互斥量而不劣化目标，最优解中必存在无同时充放电的调度策略。这一结论也由求解后逐时段诊断得到验证（第~\ref{subsec:verify}节：同时充放电时段数为 0）。

\subsubsection{模型特性}

式\eqref{eq:obj}--\eqref{eq:nonneg}构成一个线性规划问题：576 个连续变量、144 个电力平衡约束、144 个状态转移方程、288 个容量上下界、576 个充放电功率界限、144 个购电非负约束与 2 个边界条件。线性规划是凸优化问题，其局部最优解即全局最优解，无需迭代调参，求解结果可复现、可验证。

\subsection{模型求解}

模型采用 Python 语言 PuLP 建模库构建，调用开源求解器 HiGHS 求解。求解状态为 Optimal，目标值有限，求解耗时约 $0.01$ 秒，得到全局最优的日前购电策略。所得策略满足全部约束（第~\ref{subsec:verify}节给出逐项验证），完整 144 时段的购电量与分时段充放电量已按附件5 模板写入 result1.xlsx。

\subsection{购电与充放电计划}

表~\ref{tab:q1_table1}给出题目要求的 6 个指定时段的购电量及全天汇总，表~\ref{tab:q1_table2}给出 6 个四小时时段的充放电量与边界储电量。

\begin{table}[H]
  \centering
  \caption{微网在指定时间段的购电量及全天的购电量和购电费}
  \label{tab:q1_table1}
  \begin{tabular}{ccccc}
    \toprule
    时间段 & 购电量(kWh) & 时间段 & 购电量(kWh) \\
    \midrule
    10:00--10:10 & 0.00 & 12:00--12:10 & 480.41 \\
    14:00--14:10 & 0.00 & 16:00--16:10 & 445.43 \\
    18:00--18:10 & 531.89 & 20:00--20:10 & 0.00 \\
    \midrule
    \multicolumn{2}{c}{全天购电量(kWh)} & \multicolumn{2}{c}{全天购电费(元)} \\
    \multicolumn{2}{c}{59,482.70} & \multicolumn{2}{c}{35,126.95} \\
    \bottomrule
  \end{tabular}
\end{table}

\begin{table}[H]
  \centering
  \caption{储能设备在指定时间段的充放电量及 0:00 和 24:00 的储电量}
  \label{tab:q1_table2}
  \small
  \begin{tabular}{cccccc}
    \toprule
    时间段 & 充电量(kWh) & 放电量(kWh) & 时间段 & 充电量(kWh) & 放电量(kWh) \\
    \midrule
    0:00--4:00 & 4,500.00 & 0.00 & 4:00--8:00 & 833.33 & 7,073.16 \\
    8:00--12:00 & 4,787.96 & 1,892.22 & 12:00--16:00 & 5,286.04 & 101.22 \\
    16:00--20:00 & 0.00 & 6,422.37 & 20:00--24:00 & 5,333.33 & 3,177.63 \\
    \midrule
    \multicolumn{3}{c}{0:00 储电量(kWh)} & \multicolumn{3}{c}{24:00 储电量(kWh)} \\
    \multicolumn{3}{c}{6,000.00} & \multicolumn{3}{c}{6,000.00} \\
    \bottomrule
  \end{tabular}
\end{table}

表~\ref{tab:q1_table1}显示，10:00--10:10、14:00--14:10 与 20:00--20:10 三个时段购电量为零，供电完全由光伏与储能放电承担；其余三个时段购电量分别为 480.41、445.43 与 531.89 kWh。全天累计购电 $59{,}482.70$ kWh，购电费用 $35{,}126.95$ 元。表~\ref{tab:q1_table2}显示，储能全天充电集中在 0:00--4:00、8:00--16:00 与 20:00--24:00 等电价相对较低的时段，放电集中在 4:00--8:00 与 16:00--20:00 两个用电高峰时段；0:00 与 24:00 储电量均为 $6{,}000.00$ kWh，满足题目要求的边界条件。

\subsection{储能调度策略分析}

图~\ref{fig:q1_strategy}给出最优购电量与电价曲线的关系。模型呈现出显著的“谷电多购、峰电少购”特征：购电量曲线与阶梯电价曲线在时间上呈明显负相关，电价低谷时段（0:00--4:00、20:00--24:00）购电量达到全天峰值，电价高峰时段购电量大幅削减甚至为零，验证了削峰填谷的优化策略。

\begin{figure}[H]
  \centering
  \includegraphics[width=0.85\textwidth]{figures/fig_q1_1_strategy.png}
  \caption{最优购电策略与电价关系}
  \label{fig:q1_strategy}
\end{figure}

图~\ref{fig:soc}展示了储能电量与充放电功率的协同关系。储能电量全程严格运行在 $1200$--$10800$ kWh 的安全区间内，充放电功率均未越过 $\pm 5000$ kW 的物理上限，充放电行为与电量升降完全对应：0:00--4:00 夜间低电价时段集中充电使储电量快速上升；4:00--8:00 早高峰集中放电、电量明显回落并接近运行下限；白天时段电量围绕高荷电水平波动，多次触及上限 $10800$ kWh；16:00--20:00 晚高峰再次集中放电；20:00--24:00 恢复充电并于 24:00 回到初始值 $6000$ kWh。全天储电量共 21 个时段触及上限、8 个时段触及下限，表明储能容量在峰谷时段得到较充分利用，且未出现越限。

\begin{figure}[H]
  \centering
  \includegraphics[width=0.88\textwidth]{figures/fig_q1_2_soc_power.png}
  \caption{储能电量与充放电功率的协同关系}
  \label{fig:soc}
\end{figure}

\subsection{储能套利效益分析}

为量化储能的贡献，构造不使用储能的基线策略：微网不配置储能时，各时段购电量即为扣除光伏后的净负载需求，其全天购电费用为

\begin{equation}
  Z_{\mathrm{base}} = \sum_{t=1}^{144} Price(t)\cdot\max\!\big(0,\, Load(t)\Delta t - PV(t)\Delta t \big)
  \label{eq:baseline}
\end{equation}

代入附件1 数据计算得基线成本为 $48{,}052.05$ 元。有储能的最优策略购电费用为 $35{,}126.95$ 元，净节省 $12{,}925.10$ 元，节省比例达 $26.9\%$（图~\ref{fig:cost}）。储能带来的效益来自两方面：一是将低价时段电量搬运至高价时段使用（套利收益），二是将白天光伏富余电量储存供高峰使用、减少高价购电（消纳收益）。

\begin{figure}[H]
  \centering
  \includegraphics[width=0.85\textwidth]{figures/fig_q1_4_econ.png}
  \caption{有/无储能的购电成本对比}
  \label{fig:cost}
\end{figure}

从购电成本结构看（表~\ref{tab:cost_struct}），电价 $0.5$--$1.0$ 元/kWh 的中价时段承担了 $78.8\%$ 的购电量与 $86.5\%$ 的购电费用，是成本的主要来源；电价不低于 $1.0$ 元/kWh 的高价时段购电量仅占 $16.2\%$、成本占比 $10.0\%$，说明储能有效压制了高价时段的购电需求；低价时段的购电仅占 $5.0\%$，对应储能集中充电。购电量占比与成本占比的错位直观反映了套利调度的效果：购电费用向低价时段倾斜，高价时段需求被储能放电替代。

\begin{table}[H]
  \centering
  \caption{购电成本结构}
  \label{tab:cost_struct}
  \begin{tabular}{ccccc}
    \toprule
    电价区间 & 购电量(kWh) & 购电量占比 & 成本(元) & 成本占比 \\
    \midrule
    $<0.5$ 元/kWh & 2,968.05 & 5.0\% & 1,241.86 & 3.5\% \\
    $0.5$--$1.0$ 元/kWh & 46,893.86 & 78.8\% & 30,395.67 & 86.5\% \\
    $\ge 1.0$ 元/kWh & 9,620.78 & 16.2\% & 3,489.43 & 10.0\% \\
    \midrule
    合计 & 59,482.70 & 100\% & 35,126.95 & 100\% \\
    \bottomrule
  \end{tabular}
\end{table}

储能运行强度方面，全天名义总充电量为 $20{,}740.67$ kWh，总放电量为 $18{,}666.60$ kWh，充电损耗 $2{,}074.07$ kWh、放电损耗 $1{,}866.66$ kWh，总转换损耗 $3{,}940.73$ kWh；按可用容量 $9600$ kWh 折算的等效充电循环为 $2.16$ 次/天，处于合理运行强度范围。

\subsection{模型验证}
\label{subsec:verify}

为确认求解结果的正确性，对最优解执行分级验证。

\textbf{输入级验证}：附件1 数据 144 行完整，字段齐全，时间序列自 0:00 起严格每 10 分钟一点且顺序正确，无缺失、无穷或非法数值。

\textbf{求解级验证}：HiGHS 返回 Optimal 状态，目标函数值 $35{,}126.95$ 元有限。

\textbf{约束级验证}：对 144 个时段逐项复核电力平衡、状态转移、容量上下界、充放电功率界限、购电非负与终止电量边界，全部硬约束通过，失败项数为 0。

\textbf{能量守恒验证}：按“输入能量 $=$ 输出能量”核对全天能量收支，结果如表~\ref{tab:energy}所示，误差为 $0.0000$ kWh。时序层面，供给（购电、光伏与储能放电之和）与需求（负载与储能充电之和）在每个时段均严格吻合，微网在满足小区负载的同时实现了能量供需的严格平衡。

\begin{table}[H]
  \centering
  \caption{全天能量平衡验证}
  \label{tab:energy}
  \begin{tabular}{lcc}
    \toprule
    \multicolumn{2}{c}{输入能量(kWh)} & \multicolumn{1}{c}{输出能量(kWh)} \\
    \cmidrule(r){1-2}\cmidrule(l){3-3}
    购电总量 & 59,482.70 & -- \\
    光伏发电总量 & 55,482.84 & -- \\
    初始储能电量 & 6,000.00 & -- \\
    \midrule
    输入合计 & 120,965.53 & -- \\
    \midrule
    -- & -- & 负载用电总量：111,024.81 \\
    -- & -- & 终止储能电量：6,000.00 \\
    -- & -- & 充电损耗：2,074.07 \\
    -- & -- & 放电损耗：1,866.66 \\
    \midrule
    -- & -- & 输出合计：120,965.53 \\
    \bottomrule
  \end{tabular}
\end{table}

\textbf{诊断级检查}：同时充放电时段数为 0，与理论论证一致；全天供需富余量（对应弃光）仅 $0.0022$ kWh，表明光伏与储能调度高度贴合需求，几乎不存在能量浪费。

综上，问题一的最优策略在满足全部约束的前提下实现购电费用最小，求解结果正确可靠。

% 问题二：考虑紧急购电的日前购电优化模型
\section{问题二：考虑紧急购电的日前购电优化模型}

\subsection{问题二的分析}

问题二将问题一的确定性环境推广为现实中的不确定环境：每天的电价仍相同（复用附件1 的 144 点电价），但小区负载与光伏发电功率逐日变化，且 0:00 制定计划时当天的负载与光伏真值尚未可知，只能依据历史数据估计。同时，题目引入紧急购电机制——若微网实际供电低于负载，缺口须在实时阶段向外网紧急购电，电价为交易时刻电价的 5 倍；除紧急购电外的购电费用均按计划购电量结算。

与问题一相比，问题二有三处本质变化。负载与光伏由确定性预测变为逐日波动的随机量（附件2 提供 2025.2.1--12.31 共 334 天的真实数据），0:00 制定的计划购电量在信息不完整的情况下必须锁定全天；供需缺口不再作为硬约束，而转化为 5 倍电价的高额惩罚，目标函数中新增紧急购电费用项；储能的角色也从“确定性套利”升级为“对冲不确定性”——实际负载高于预期时放电减少紧急购电，光伏富余时吸收电量避免浪费。由于缺口惩罚（5 倍）远高于正常购电（1 倍），计划购电量与储能调度的权衡呈现明显的非对称结构，属于典型的日前--实时两阶段随机决策问题\cite{birge1997stochastic,hou2023multi}。

据此，本文将问题二分解为两个耦合层次：\textbf{日前决策层}在每天 0:00 基于历史情景制定计划购电量 $E_{plan}(t)$；\textbf{日内执行层}在真实负载、光伏逐时段揭晓后，按因果规则实时调度储能充放电，缺口由紧急购电补齐。两层之间通过“计划购电量锁定、储能实时可调”的接口衔接。

\subsection{数据准备}

附件2 提供 2025.1.1--12.31 全年 365 天、每天 144 个 10 分钟时段的小区负载与光伏发电实际功率。按题目要求，规划区间为 2025.2.1--12.31，共 334 天。数据审查表明附件2 无缺失、无负值，可直接使用。

对规划区间内净负载（负载与光伏之差）的统计显示，净负载日总量均值为 $55{,}542$ kWh、标准差为 $23{,}967$ kWh（约为均值的 $43\%$），日际波动显著，这是问题二不确定性的主要来源。进一步考察日前可估计性：仅利用历史数据对光伏的归一化均方根误差（nRMSE）约为 $0.034$，高度可估计；负载的 nRMSE 约为 $0.339$，随机波动主要来自负载。此外，光伏日发电量呈明显的季节性（1 月与 12 月约 $3.6\times10^4$ kWh，6 月与 9 月约 $6\times10^4$ kWh），且 $19.3\%$ 的时段光伏出力超过负载，白天富余电量可被储能吸收。

\subsection{日前决策层：历史联合残差情景的随机优化模型}

\subsubsection{情景构造}

0:00 制定计划时，当天负载与光伏的真值未知，需要用历史数据刻画其可能分布。本文采用\textbf{历史联合残差}法构造情景：先分别用历史数据生成当天预测，再以前 30 天的历史预测误差作为扰动，生成 $M=30$ 个等权情景（$\pi_s=1/M$）。具体步骤如下：

\textbf{步骤 1（历史预测）}：对规划日 $d$，负载预测取前 35 天内同星期日对应时段的均值，光伏预测取前 7 天对应时段的均值；对每个历史日 $i<d$，用 $i$ 日当时可获得的信息同样生成其历史预测 $\hat{L}_i(t)$、$\hat{V}_i(t)$，保证不泄露未来信息。

\textbf{步骤 2（历史联合残差）}：对每个历史日 $i$ 记录其 144 点预测误差序列
\begin{equation}
  \varepsilon_i^{L}(t) = L_i(t) - \hat{L}_i(t),\qquad
  \varepsilon_i^{V}(t) = V_i(t) - \hat{V}_i(t),\quad t=1,\ldots,144
  \label{eq:residual}
\end{equation}
负载误差与光伏误差取自同一历史日，保留了两者的联合相关结构（如阴雨天气对二者的同步影响）。

\textbf{步骤 3（当天情景生成）}：对规划日 $d$，先生成当天预测 $\hat{L}_d(t)$、$\hat{V}_d(t)$，再以历史残差构造情景：
\begin{equation}
  L_d^{s}(t) = \big[\hat{L}_d(t) + \varepsilon_{i(s)}^{L}(t)\big]_{+},\qquad
  V_d^{s}(t) = \big[\hat{V}_d(t) + \varepsilon_{i(s)}^{V}(t)\big]_{+}
  \label{eq:scenario}
\end{equation}
其中 $i(s)$ 为情景 $s$ 对应的历史日（取 $d$ 前 30 天），$[\cdot]_{+}$ 表示负值截断为零。

该情景构造方法数据驱动、无需假设误差分布形式，同时保留了日内时序相关性与负载--光伏联合误差结构，与两阶段随机优化中基于历史样本的经验分布建模思想一致\cite{garciatorres2021stochastic}。

\subsubsection{决策变量}

日前层的唯一核心决策变量是计划购电量 $E_{plan}(t)\ge 0$（$t=1,\ldots,144$），0:00 锁定、全天不可更改，这是题目的明确规定。为快速评估不同计划的期望成本，为每个情景 $s$ 引入仅用于评估的第二阶段变量：储能充电量 $E_{charge}^{s}(t)$、放电量 $E_{discharge}^{s}(t)$、储电量 $E_{storage}^{s}(t)$ 与紧急购电量 $E_{emg}^{s}(t)$。这些情景相关变量不代表实际执行，实际执行由日内层的动态规划执行器按因果规则完成（见第~\ref{subsec:dp}节）。

\subsubsection{目标函数}

日前层在计划购电费用与期望紧急购电费用之间权衡：

\begin{equation}
  \min_{E_{plan},\{E^{s}\}}\;
  \sum_{t=1}^{144} Price(t)\,E_{plan}(t)
  + \sum_{s=1}^{M}\pi_s \sum_{t=1}^{144} \kappa\, Price(t)\, E_{emg}^{s}(t)
  \label{eq:q2_obj}
\end{equation}

其中 $\kappa=5$ 为紧急购电电价倍数，$\pi_s=1/M$。紧急购电费用的系数为正常购电的 5 倍，模型因此倾向于通过适度增加计划购电量与充分利用储能来压低高额惩罚。

\subsubsection{约束条件}

对每个情景 $s$ 与每个时段 $t$：

\textbf{(1) 电力平衡（含紧急购电，允许弃光）}：
\begin{equation}
  E_{plan}(t) + E_{emg}^{s}(t) + PV_d^{s}(t)\Delta t + \eta E_{discharge}^{s}(t)
  \ge Load_d^{s}(t)\Delta t + E_{charge}^{s}(t)
  \label{eq:q2_balance}
\end{equation}

\textbf{(2) 储能状态转移方程}（继承问题一的效率口径）：
\begin{equation}
  E_{storage}^{s}(t) = E_{storage}^{s}(t-1) + \eta E_{charge}^{s}(t) - E_{discharge}^{s}(t),\quad
  E_{storage}^{s}(0) = 6000
  \label{eq:q2_state}
\end{equation}

\textbf{(3) 容量与功率约束}（继承问题一）：
\begin{equation}
  1200 \le E_{storage}^{s}(t) \le 10800,\qquad
  0 \le E_{charge}^{s}(t),\, E_{discharge}^{s}(t) \le 833.33
  \label{eq:q2_bounds}
\end{equation}

\textbf{(4) 日末电量约束}：
\begin{equation}
  E_{storage}^{s}(144) = 6000,\quad \forall s
  \label{eq:q2_terminal}
\end{equation}

\textbf{(5) 非负与非提前性}：$E_{plan}(t)\ge 0$，$E_{emg}^{s}(t)\ge 0$；计划购电量 $E_{plan}(t)$ 对全部情景取同一值（第一阶段变量天然满足非提前性）。

式\eqref{eq:q2_obj}--\eqref{eq:q2_terminal}构成一个两阶段随机线性规划，第一、二阶段变量共 $144+30\times576=17{,}424$ 个，规模适中，HiGHS 求解器可在秒级求得全局最优解。

\subsection{日内执行层：动态规划价值执行器}
\label{subsec:dp}

\subsubsection{实时决策问题}

日内执行阶段，每个时段 $t$ 的真实负载与光伏逐次揭晓。定义该时段的净缺口
\begin{equation}
  r_t = Load_t\,\Delta t - PV_t\,\Delta t - E_{plan}(t)
  \label{eq:netgap}
\end{equation}
$r_t>0$ 表示计划供给不足，需要储能放电或紧急购电补齐；$r_t<0$ 表示供给富余，可向储能充电。执行器的任务是在因果信息约束下决定充电量 $C_t$ 或放电量 $D_t$（二者互斥），最小化当天的紧急购电费用。

\subsubsection{历史路径价值函数}

储能调度的难点在于当前决策的跨时段影响：当前放电缓解本时段缺口，却降低未来时段的可用电量。沿用问题一的库存价值思想，本文用历史情景构造储能电量的\textbf{未来价值函数}。对每条历史情景 $\omega$，自日末向前递推：
\begin{equation}
  H_{t,\omega}(e) = \min_{x\in A_t(e,r_{t,\omega})}\Big\{
    5\,Price(t)\,\big[r_{t,\omega}+\psi(x)\big]_{+} + H_{t+1,\omega}(e+x)\Big\}
  \label{eq:dp_recur}
\end{equation}
其中 $e$ 为当前储电量，$x$ 为储能内部电量增量（充电为正、放电为负），$A_t(e,r)$ 为满足容量与功率约束且单向充放电的可行增量集合，$\psi(x)$ 为效率转换函数（充电 $x/\eta$、放电 $\eta x$）。终端价值取线性形式
\begin{equation}
  H_{T+1,\omega}(e) = -v\,e,\qquad v = 0.48\ \text{元/kWh}
  \label{eq:terminal_value}
\end{equation}
即日末每 1 kWh 库存的跨日边际价值约为 0.48 元。每条路径的价值函数 $H_{t,\omega}(e)$ 为凸分段线性函数且随 $e$ 单调非增（库存越多，未来补缺费用越低），对全部历史情景平均得到
\begin{equation}
  \bar{H}_t(e) = \frac{1}{M}\sum_{\omega=1}^{M} H_{t,\omega}(e)
  \label{eq:avg_value}
\end{equation}
凸性与单调性在平均下保持，$\bar{H}_t(e)$ 以折线形式存储。

\subsubsection{实时执行}

进入实际时段 $t$，观测到真实净缺口 $r_t$ 与当前储电量 $e=SOC_{t-1}$ 后，求解单时段优化
\begin{equation}
  x_t^{*} \in \arg\min_{x\in A_t(e,r_t)}\Big\{
    5\,Price(t)\,\big[r_t+\psi(x)\big]_{+} + \bar{H}_{t+1}(e+x)\Big\}
  \label{eq:dp_exec}
\end{equation}
并恢复储能动作 $C_t=[x_t^{*}]_{+}/\eta$、$D_t=[-x_t^{*}]_{+}\cdot\eta$，实际紧急购电量为 $E_{emg}(t)=[r_t+\psi(x_t^{*})]_{+}$。式\eqref{eq:dp_exec}中当前代价仅依赖当时段观测，未来价值 $\bar{H}_{t+1}$ 于每日 0:00 由历史情景构建、不依赖当天未来真值，故执行器严格满足因果性（非提前性）。直观上，解析意义是：当前放电既减少本时段 5 倍电价的紧急购电，也损耗未来库存的期望价值，DP 执行器在二者间做动态权衡。

\subsubsection{对照执行器}

为隔离日内执行层的贡献，本文实现两个对照执行器，与 DP 执行器共用同一日前计划：\textbf{解析响应}执行器贪婪地最小化当前紧急购电——富余时段尽量充电、缺口时段尽量放电，不考虑未来价值；\textbf{MPC} 执行器利用历史残差估计相邻时段相关性，对剩余时段净负荷做衰减修正后滚动优化。三者在相同日前计划下回放真实数据，即可比较执行策略本身的优劣。

\subsection{跨日边界条件与全年回测框架}

问题一要求 0:00 与 24:00 储电量相同；问题二未明确每日边界，若简单沿用“每日强制回到 6000 kWh”的独立日假设，将迫使储能每天回到初始状态，损失跨日套利机会。本文采用\textbf{跨日连续 + 年度循环}的边界处理：首日以 6000 kWh 起始，相邻两日储电量物理衔接（当日 24:00 即次日 0:00），保持全年运行连续：

\begin{equation}
  E(1,0) = 6000,\qquad
  E(d,0) = E(d-1,144),\ d=2,\ldots,334
  \label{eq:year_boundary}
\end{equation}

年度循环性则由日内层的终端价值项\eqref{eq:terminal_value}软引导实现：末日的储能决策在“当前紧急购电代价”与“终端库存跨年价值 $v e$”之间权衡，使年末储电量自然回到年初水平附近（实际年末储电量 6069.60 kWh，相对 6000 kWh 目标偏差 $1.16\%$）。

在此框架下对 334 天进行完整回测：每日 0:00 求解日前随机 LP 得 $E_{plan}$，构建历史路径价值函数 $\bar{H}_t$，随后按真实数据逐时段执行 DP 策略并记录实际充放电与紧急购电；各天结果按跨日约束衔接，汇总得到全年指标。

\subsection{模型求解}

模型采用 Python 语言实现：日前层用 PuLP 建模、HiGHS 求解（334 个两阶段 LP，单日规模约 $1.7\times10^4$ 变量，均返回 Optimal）；日内层用 NumPy 实现 DP 折线递推与实时决策。求解全部完成后，按附件5 模板将 334 天的计划购电量（144 时段）、分时段充放电量与紧急购电记录写入 result2.xlsx。

\subsection{结果分析}

\subsubsection{年度购电成本结构}

334 天回测的年度总购电费用为 \textbf{14,120,324 元}，其中计划购电费用 12,870,635 元（占 $91.1\%$），紧急购电费用 1,249,689 元（占 $8.9\%$）。成本分解如图~\ref{fig:q2_cost}所示：紧急购电费用占比约一成，说明两阶段策略在多数时段通过储能调节成功避免了高额惩罚，仅在少数缺口较大时段触发紧急购电。

\begin{figure}[H]
  \centering
  \includegraphics[width=0.92\textwidth]{figures/fig_q2_1_cost.png}
  \caption{问题二年度购电成本分解}
  \label{fig:q2_cost}
\end{figure}

\subsubsection{储能的年经济价值}

为量化储能在不确定性环境下的贡献，构造无储能基线：微网不配置储能时，实际缺口全部由紧急购电承担。计算得无储能年度成本为 16,407,320 元，而两阶段策略为 14,120,324 元，全年净节省 \textbf{2,286,996 元}，节省比例达 \textbf{13.94\%}。与问题一单日 $26.9\%$ 的节省比例相比，问题二中储能的收益部分被预测偏差所侵蚀，但年化节省仍超过两百万元，表明储能套利与不确定性对冲的组合价值显著。

\subsubsection{指定日期的购电策略}

表~\ref{tab:q2_days}、表~\ref{tab:q2_days2}与表~\ref{tab:q2_days3}分别按题目表1、表2、表3 的格式汇总四个指定日期（2025.3.20、2025.6.21、2025.9.23、2025.12.21）的购电量、充放电量与紧急购电量。

\begin{table}[H]
  \centering
  \caption{指定日期的计划购电量（表1 格式）}
  \label{tab:q2_days}
  \small
  \begin{tabular}{lcccc}
    \toprule
    时间段 & 2025.3.20 & 2025.6.21 & 2025.9.23 & 2025.12.21 \\
    \midrule
    10:00--10:10  & 0.00 & 0.00 & 0.00 & 0.00 \\
    12:00--12:10  & 577.97 & 0.00 & 275.35 & 924.77 \\
    14:00--14:10  & 0.00 & 0.00 & 0.00 & 0.00 \\
    16:00--16:10  & 454.39 & 121.97 & 470.86 & 835.18 \\
    18:00--18:10  & 701.56 & 436.85 & 785.13 & 735.86 \\
    20:00--20:10  & 21.06 & 0.00 & 151.90 & 0.00 \\
    \midrule
    全天购电量(kWh) & 66,188.41 & 35,517.60 & 64,271.75 & 96,508.07 \\
    全天购电费(元) & 41,848.06 & 21,064.80 & 44,734.10 & 62,467.69 \\
    \bottomrule
  \end{tabular}
\end{table}

\begin{table}[H]
  \centering
  \caption{指定日期的储能充放电量与边界储电量（表2 格式）}
  \label{tab:q2_days2}
  \footnotesize
  \begin{tabular}{lcccccccc}
    \toprule
    \multirow{2}{*}{时间段} & \multicolumn{2}{c}{2025.3.20} & \multicolumn{2}{c}{2025.6.21} & \multicolumn{2}{c}{2025.9.23} & \multicolumn{2}{c}{2025.12.21} \\
    \cmidrule(lr){2-3}\cmidrule(lr){4-5}\cmidrule(lr){6-7}\cmidrule(lr){8-9}
    & 充电量 & 放电量 & 充电量 & 放电量 & 充电量 & 放电量 & 充电量 & 放电量 \\
    \midrule
    0:00--4:00   & 8,007.65 & 183.56 & 816.19 & 58.87 & 7,098.02 & 39.44 & 8,982.65 & 119.09 \\
    4:00--8:00   & 1,873.33 & 4,806.84 & 1,388.01 & 1,725.67 & 1,742.29 & 5,655.40 & 2,509.92 & 1,484.03 \\
    8:00--12:00  & 6,049.91 & 2,777.38 & 8,818.73 & 0.00 & 3,033.55 & 1,896.62 & 5,244.87 & 6,461.75 \\
    12:00--16:00 & 4,274.33 & 352.57 & 0.00 & 0.00 & 5,499.14 & 501.66 & 6,338.99 & 2,234.15 \\
    16:00--20:00 & 84.78 & 5,848.65 & 201.01 & 5,035.22 & 353.45 & 5,102.80 & 14.40 & 5,790.85 \\
    20:00--24:00 & 221.37 & 2,941.48 & 944.09 & 2,515.96 & 599.16 & 2,538.41 & 333.07 & 2,806.15 \\
    \midrule
    0:00 储电量   & \multicolumn{2}{c}{1,529.19} & \multicolumn{2}{c}{2,862.18} & \multicolumn{2}{c}{2,675.92} & \multicolumn{2}{c}{1,474.55} \\
    24:00 储电量  & \multicolumn{2}{c}{1,200.00} & \multicolumn{2}{c}{3,440.40} & \multicolumn{2}{c}{1,686.40} & \multicolumn{2}{c}{1,560.50} \\
    \bottomrule
  \end{tabular}
\end{table}

\begin{table}[H]
  \centering
  \caption{指定日期的紧急购电量（表3 格式）}
  \label{tab:q2_days3}
  \small
  \begin{tabular}{lcc}
    \toprule
    日期 & 紧急购电时间段 & 紧急购电量(kWh) \\
    \midrule
    2025.3.20  & 20:30--20:40 & 132.80 \\
    2025.6.21  & -- & 0.00 \\
    2025.9.23  & 20:20--20:30 & 221.54 \\
    2025.12.21 & -- & 0.00 \\
    \bottomrule
  \end{tabular}
\end{table}

\subsubsection{紧急购电的时空特征}

全年共 158 天发生紧急购电（占总天数 $47.3\%$），总紧急购电量 246,579 kWh，约占全年总负荷的 $1.23\%$。日度购电费用的分布高度右偏：全年日均购电费用 42,276 元、标准差 16,653 元（变异系数 $39.4\%$），最低日费用 12,297 元（4 月 12 日）、最高日费用 112,099 元（6 月 1 日），少数极端高成本日贡献了紧急费用的大部分，对应光伏出力大幅低于预期的情形。分季节统计显示，夏季紧急购电费用最高（约 39 万元），与其光伏波动与高温负荷特征一致；冬季紧急触发频率最高（$52.5\%$），与光伏日发电量全年最低的季节特征一致。紧急购电费用与频率的季节错位说明：费用由极端缺口事件驱动，频率由常规预测偏差驱动。各季节的储能运行均遵循“谷时充电、峰时放电”的基本策略，但幅度存在差异：夏季光伏充足时段储电量上升更快、白天更多时段处于高荷电水平；冬季光伏出力有限，储能更依赖夜间低电价时段充电。

\subsubsection{储能全年运行状态}

全年储电量均值为 6,489 kWh（约为可用容量的 $60.4\%$），标准差 3,403 kWh；$31.11\%$ 的时段储电量贴近容量边界（1200 或 10800 kWh），说明储能容量得到充分利用，且全部时段均未违反容量约束。图~\ref{fig:q2_soc_year}展示全年储电量的演化轨迹，储能随季节在 1200--10800 kWh 之间往复运行，运行形态与问题一一致；年末储电量为 6069.60 kWh，与年初 6000 kWh 的循环目标偏差 $1.16\%$，由终端价值软引导实现。

\begin{figure}[H]
  \centering
  \includegraphics[width=0.92\textwidth]{figures/fig_q2_7_soc_year.png}
  \caption{全年储电量演化轨迹}
  \label{fig:q2_soc_year}
\end{figure}

\subsubsection{日前评估与实际执行的偏差}

日前 LP 评估的期望成本与实际 DP 执行成本之间存在系统性偏差：全年计划购电中有 246,579 kWh 未能在实时阶段按原定用途使用（利用率缺口约 $1.19\%$）。这一偏差源于两阶段的模型一致性限制——日前层假设储能可“看到”完整情景并最优响应，而实际执行层只能按因果规则逐时段决策，故日前层对储能调节能力的评估偏乐观。$1.19\%$ 的计划浪费率在可接受范围内，且已通过紧急购电机制在实时阶段完成闭环。

\subsection{模型验证}
\label{subsec:verify2}

\textbf{跨日连续性验证。} 对 334 天逐日核验式\eqref{eq:year_boundary}：333 个跨日过渡点的储电量衔接误差均小于 $10^{-6}$ kWh，物理连续性完美满足。

\textbf{硬约束验证。} 对全部 $334\times144$ 个时段逐项复核电力平衡、储能状态转移、容量界限与功率界限，违反次数为 0；储能全年充放电循环效率为 $81.00\%$，与理论往返效率 $\eta^2=0.81$ 完全吻合，验证了状态方程与效率口径的物理正确性。

\textbf{执行器对照实验。} 固定同一日前计划，分别用解析响应、MPC 与 DP 价值执行器回放全年真实数据。三种执行器的年度总成本分别为 14,174,236 元、14,278,883 元与 14,022,232 元：DP 执行器相比解析响应改进 \textbf{1.07\%}，相比 MPC 改进 \textbf{1.80\%}；与不使用动态执行器的固定计划基准（约 16,364,000 元）相比改进 \textbf{14.31\%}，配对检验 $p<0.0001$，效应量 Cohen's $d=2.704$（大效应）。这表明 DP 执行器对“当前紧急代价与未来库存价值”的动态权衡显著优于只考虑当前代价的贪婪策略，日内执行层的独立贡献得到统计验证。

\textbf{边界条件对比。} 为考察边界条件的影响，本文比较了三种设定下的年度总成本：每日强制回 6000 kWh（V1）为 14,249,906 元；完全无终端约束（V2）为 14,118,116 元，但年末储电量贴至下限 1312.01 kWh，年度运行不可循环；跨日连续 + 年度循环（V3）为 14,120,324 元，相比 V1 降低 $0.91\%$，相比 V2 仅多 $0.016\%$，却保证了年度可循环性。V3 以几乎可忽略的成本代价兼顾了经济性与运行可行性，作为最终方案。

\subsection{模型评价}

问题二的模型具有清晰的层次结构与严格的因果性保证：日前层以历史联合残差情景刻画不确定性，两阶段 LP 求解高效、全局最优；日内层 DP 执行器只使用当前观测与历史构建的未来价值函数，无未来信息泄露。全年回测与四级验证（输入、求解、硬约束、诊断）确认结果正确可信。

模型的局限同样明确。日前评估层与实际执行层存在模型不一致（计划浪费率 $1.19\%$），回测已验证偏差可接受，理论上仍可通过迭代校准进一步缩小。此外，紧急购电费用占比 $8.9\%$ 高于采用分位数准则的报童式方法的理论水平（约 $5.77\%$），差距源于本文采用期望最小化准则——期望准则对非对称惩罚的刻画不如分位数准则，这是方法选择层面的已知权衡。

% 问题三：多时点预报与滚动调整购电模型
\section{问题三：多时点预报下的滚动调整购电模型}

\subsection{问题三的分析}

问题三的信息结构比前两问更为丰富：除 0:00 的日前预报外，每天 6:00、12:00 与 18:00 均可获得未来 24 小时的光伏功率预报（附件3），微网可根据新预报调整剩余时段的购电策略。调整伴随成本：计划购电量高于调整购电量的部分收取交易时刻电价 $50\%$ 的违约费用，低于调整购电量的部分按 $1.5$ 倍电价计费；总费用为计划购电、紧急购电与调整费用之和。题目同时要求分析是否需要在 0:00 以外的时刻制定调整策略。

问题三的本质是一个多阶段动态决策问题：预报信息随时间逐步更新，每个决策时刻都面临“维持原计划”与“按新信息调整”的权衡——调整可降低供需偏差与紧急购电费用，但需支付违约或加价费用。本文将问题三建模为四阶段确定性滚动优化：各决策时刻以当时可获得的最新预报重新求解剩余时段的购电优化问题，仅当预报改善度超过阈值时执行调整，以控制频繁重优化的计算与违约成本。

\subsection{多阶段滚动调整购电模型}

\subsubsection{决策框架}

每天设置四个决策时刻 $k\in\{0,6,12,18\}$（以小时计，$k=0$ 即 0:00）。各时刻的决策规则为：

\textbf{(1) 0:00 日前计划}：基于 0:00 的光伏预报与历史负载估计，以“计划购电费 $+$ 期望紧急购电费”最小为目标求解问题二式两阶段线性规划，得到全天计划购电量 $E_{plan}(t)$。

\textbf{(2) 6:00/12:00/18:00 滚动调整}：设当前时刻 $k$ 前的最新执行购电量为 $E_{base}(t)$，当新预报相对上一预报的\textbf{平均光伏功率改善度}超过阈值 $\theta$（本文取 $\theta=1000$ kW）时触发重优化：以剩余时段 $t\ge t_k$ 为决策区间重新求解购电优化，得到调整购电量 $E_{adjust}(t)$；否则维持原计划。触发规则以计算效率为目标，通过参数扫描确定阈值。

\subsubsection{目标函数与调整成本}

全天总费用由三部分构成：

\begin{equation}
  Z_{\mathrm{total}} = Z_{\mathrm{plan}} + Z_{\mathrm{adjust}} + Z_{\mathrm{emergency}}
  \label{eq:q3_obj}
\end{equation}

其中计划购电费按 0:00 计划结算，紧急购电费为实时缺口按 $5$ 倍电价的补充费用（同问题二），调整费用为：

\begin{equation}
  Z_{\mathrm{adjust}} = \sum_{k\in\{6,12,18\}}\sum_{t=t_k}^{144}\Big[
    1.5\, Price(t)\,\delta_k^{+}(t) - 0.5\, Price(t)\,\delta_k^{-}(t)\Big]
  \label{eq:q3_adjust}
\end{equation}

式中 $\delta_k^{+}(t)=\max(0,\,E_{adjust,k}(t)-E_{base,k}(t))$ 为增购量（按 $1.5$ 倍电价计费），$\delta_k^{-}(t)=\max(0,\,E_{base,k}(t)-E_{adjust,k}(t))$ 为减购量（退还 $50\%$ 电价）。储能约束、效率口径（$\eta=0.9$）与电力平衡结构均继承问题二。

\subsubsection{实际执行的购电量与费用结算}

实际执行以各时刻最新策略为准：时段 $t$ 的最终购电量 $E_{final}(t)$ 取该时段所属决策区间的最新调整值；日内实时供需缺口按真实负载与光伏结算，不足部分触发紧急购电：

\begin{equation}
  E_{emergency}(t) = \max\Big(0,\;
  Load_{real}(t)\Delta t + E_{charge}(t) - E_{final}(t) - PV_{real}(t)\Delta t - \eta E_{discharge}(t)\Big)
  \label{eq:q3_emg}
\end{equation}

储能充放电由继承自问题二的实时调度规则执行。

\subsection{求解与验证}

按上述框架对 2025.2.1--12.31 共 334 天滚动求解，全部天数获得最优解，平均单日求解时间 $8.4$ 秒，总计约 $46.8$ 分钟。模型验证覆盖三个方面。逐时段硬约束复核（储能电量、充放电功率）全部满足；储电量跨时段一致性复核误差小于 $0.01$ kWh。能量守恒审计进一步表明，将光伏富余未消纳部分显式计入弃光项后，全年能量平衡完全闭合：

\begin{equation}
  E_{grid} + E_{PV} + \eta E_{discharge} = E_{load} + E_{charge} + E_{curtail}
  \label{eq:q3_balance}
\end{equation}

全年弃光总量为 $2{,}309{,}542$ kWh（占光伏总发电的 $12.11\%$），日均弃光 $6{,}915$ kWh（中位数 $4{,}357$），$66.8\%$ 的日期弃光率超过 $5\%$，弃光集中于光伏高峰时段。

\subsection{结果分析}

\subsubsection{年度成本结构}

全年总购电费用为 \textbf{16,489,545 元}，其中计划购电费用 14,968,290 元（$90.8\%$）、紧急购电费用 1,214,641 元（$7.4\%$）、调整费用 306,613 元（$1.9\%$），日均成本 49,370 元（图~\ref{fig:q3_cost}）。调整费用占比不足 $2\%$，表明调整机制以较小的代价参与了费用控制。

\begin{figure}[H]
  \centering
  \includegraphics[width=0.62\textwidth]{figures/fig_q3_1_cost.png}
  \caption{问题三全年购电费用构成}
  \label{fig:q3_cost}
\end{figure}

全年计划购电量 19{,}783{,}693 kWh，调整后最终购电量 20{,}012{,}913 kWh（净增 $1.16\%$）；紧急购电 446{,}763 kWh，占总购电量的 $2.23\%$。

\subsubsection{调整策略的执行特征}

全年共 137 天（$41.0\%$）触发了至少一次策略调整。以 Day 8（2025-02-08）为例：6:00 新预报显示光伏出力低于 0:00 预报且储电量偏低，触发重优化，该日增购 309 kWh，调整费用 506.72 元。调整集中在预报偏差较大的日期，且多发生在 6:00（当日首个更新时刻），与预报误差随提前量增大的一般规律一致。

\subsubsection{储能运行与全年轨迹}

储能初始电量 6000 kWh，年末降至 4002 kWh（全年净放电 1{,}998 kWh，主要由充放电效率损耗累积所致），全年运行区间为 $2006$--$7434$ kWh，位于题目规定的安全范围内。图~\ref{fig:q3_soc}给出全年日末储电量轨迹：储电量随光伏季节性呈现“春夏高、冬季低”的走势，配合全年 $12.11\%$ 的弃光率，说明在光伏高峰季节储能的消纳能力接近饱和，是弃光的主要来源之一。

\begin{figure}[H]
  \centering
  \includegraphics[width=0.9\textwidth]{figures/fig_q3_2_soc.png}
  \caption{问题三全年日末储电量轨迹}
  \label{fig:q3_soc}
\end{figure}

\subsubsection{指定日期的购电策略}

表~\ref{tab:q3_days}与表~\ref{tab:q3_days2}按题目表1、表2 格式汇总四个指定日期（2025.3.20、2025.6.21、2025.9.23、2025.12.21）的购电量与储能充放电量。

\begin{table}[H]
  \centering
  \caption{指定日期的购电量与全天购电费（表1 格式，购电量单位：kWh）}
  \label{tab:q3_days}
  \small
  \begin{tabular}{lcccc}
    \toprule
    时段 & 2025.3.20 & 2025.6.21 & 2025.9.23 & 2025.12.21 \\
    \midrule
    10:00--10:10 & 0.00 & 1.39 & 0.00 & 361.36 \\
    12:00--12:10 & 8.19 & 1.39 & 0.00 & 321.02 \\
    14:00--14:10 & 44.56 & 1.39 & 0.00 & 306.93 \\
    16:00--16:10 & 634.57 & 73.55 & 552.83 & 807.65 \\
    18:00--18:10 & 586.87 & 377.64 & 686.76 & 634.90 \\
    20:00--20:10 & 617.03 & 539.01 & 709.96 & 637.49 \\
    \midrule
    全天购电量 & 64,570.35 & 34,663.12 & 64,280.85 & 93,124.04 \\
    全天购电费(元) & 48,849.13 & 25,503.23 & 50,287.53 & 70,627.79 \\
    \bottomrule
  \end{tabular}
\end{table}

\begin{table}[H]
  \centering
  \caption{指定日期的储能充放电量与边界储电量（表2 格式，单位：kWh）}
  \label{tab:q3_days2}
  \footnotesize
  \begin{tabular}{lcccccccc}
    \toprule
    \multirow{2}{*}{时间段} & \multicolumn{2}{c}{2025.3.20} & \multicolumn{2}{c}{2025.6.21} & \multicolumn{2}{c}{2025.9.23} & \multicolumn{2}{c}{2025.12.21} \\
    \cmidrule(lr){2-3}\cmidrule(lr){4-5}\cmidrule(lr){6-7}\cmidrule(lr){8-9}
    & 充电量 & 放电量 & 充电量 & 放电量 & 充电量 & 放电量 & 充电量 & 放电量 \\
    \midrule
    0:00--4:00   & 1,729 & 0 & 28 & 1,257 & 201 & 104 & 3,111 & 0 \\
    4:00--8:00   & 542 & 1,319 & 333 & 964 & 97 & 1,424 & 1,451 & 840 \\
    8:00--12:00  & 1,389 & 1,306 & 3,333 & 0 & 1,417 & 285 & 1,097 & 326 \\
    12:00--16:00 & 2,882 & 0 & 3,035 & 28 & 2,132 & 250 & 910 & 403 \\
    16:00--20:00 & 208 & 1,903 & 111 & 1,250 & 7 & 1,743 & 7 & 2,056 \\
    20:00--24:00 & 201 & 1,368 & 35 & 2,174 & 132 & 1,438 & 347 & 979 \\
    \midrule
    0:00 储电量   & \multicolumn{2}{c}{4,005} & \multicolumn{2}{c}{4,035} & \multicolumn{2}{c}{5,935} & \multicolumn{2}{c}{4,002} \\
    24:00 储电量  & \multicolumn{2}{c}{4,365} & \multicolumn{2}{c}{4,550} & \multicolumn{2}{c}{4,279} & \multicolumn{2}{c}{5,629} \\
    \bottomrule
  \end{tabular}
\end{table}

\begin{table}[H]
  \centering
  \caption{指定日期的紧急购电汇总（表3 格式）}
  \label{tab:q3_days3}
  \footnotesize
  \begin{tabular}{lcccc}
    \toprule
    日期 & 紧急购电时段数 & 总购电量(kWh) & 总费用(元) & 最集中时段(购电量) \\
    \midrule
    2025.3.20 & 79 & 693.53 & 2,564.13 & 05:00--05:50 (109.82 kWh) \\
    2025.6.21 & 60 & 836.30 & 2,216.93 & 04:00--04:50 (334.61 kWh) \\
    2025.9.23 & 72 & 503.69 & 1,299.61 & 01:30--01:50 (151.25 kWh) \\
    2025.12.21 & 53 & 135.44 & 600.56 & 19:00--19:50 (47.06 kWh) \\
    \bottomrule
  \end{tabular}
\end{table}

四个节气的紧急购电呈现差异化的时段分布：春分集中在早晨 08:00--09:00（单小时 157.77 kWh），对应负荷快速上升而光伏尚未启动的时段；夏至集中在凌晨 04:00--05:00（334.61 kWh，单日最高），对应光伏预测在清晨时段偏差较大的情形；冬至时段数最少（53 个），但集中于晚高峰 19:00--20:00，平均紧急购电单价最高（$4.43$ 元/kWh，对应高电价时段），与春秋季的凌晨低价时段（约 $2.6$ 元/kWh）形成对照。紧急购电主要发生在负荷快速变化与光伏预测不确定性较高的时段，这一分布特征可为购电策略的进一步优化提供参考。

\subsection{调整策略的必要性分析}

题目要求回答“是否需要引入其他时刻的预报制定调整购电策略”，本文的结果支持肯定回答。全年 $41.0\%$ 的日期触发了调整，预报信息更新是高频且显著的；在调整机制运行下，紧急购电量占比被控制在 $2.23\%$（费用占比 $7.4\%$），实时供需失衡的代价处于较低水平；而调整费用仅占总费用的 $1.9\%$，引入调整策略的额外成本远低于其所缓解的供需失衡代价。据此，在 6:00、12:00 与 18:00 根据新预报滚动调整购电策略是必要且经济的。

% 问题四：波动电价下的购电策略模型
\section{问题四：波动电价下的购电策略模型}

\subsection{问题四的分析}

实际问题中，外网电价并非逐日固定，而是实时波动的。问题四要求针对波动电价重新建立购电策略模型，并在波动电价下分别重算问题二（对应结果文件 result4-2.xlsx）与问题三（对应 result4-3.xlsx）。

波动电价对购电决策的影响体现在两个层面。日前计划的购电费用不再是“购电量 $\times$ 固定电价”，而是逐时段购电量与对应时段电价乘积之和，低价时段多购、高价时段少购成为额外的优化维度；日内执行层的紧急购电惩罚基准随电价波动，储能放电的价值随电价水平动态变化，套利策略需要追踪日内价格信号。值得注意的是，波动电价并未改变问题二、问题三的决策结构与约束体系，两者唯一的本质变化是电价从固定序列 $Price$ 变为逐日逐时段的波动序列 $Price_d(t)$。因此，本文在问题二、问题三冻结模型的基础上，将电价输入替换为附件4 的波动电价数据，以最小改动完成问题四的建模：Q4-2 即“问题二模型 + 波动电价”，Q4-3 即“问题三模型 + 波动电价”。本节给出 Q4-2 的完整结果；Q4-3 的模型框架与其一致，其数值结果待问题三求解完成后一并补充。

\subsection{波动电价数据}

附件4 提供 2025.1.1--12.31 全年 365 天、每天 144 个 10 分钟时段的电价数据，共 52,560 个电价值，规划区间仍取 2025.2.1--12.31（334 天）。数据审查显示：电价全部为正，取值区间为 $[0.0076,\ 1.7936]$ 元/kWh，均值 $0.7662$ 元/kWh，峰谷比约 236 倍，变异系数 $42.3\%$，波动强度显著高于附件1 的单日分时电价（$[0.371,\ 1.395]$），且电价同时呈现日内分时特征与日间、月间波动特征。

\subsection{Q4-2：波动电价下的两阶段购电模型}

Q4-2 完全继承问题二的两阶段架构：日前层基于历史联合残差情景（$M=30$）求解两阶段随机线性规划，锁定计划购电量 $E_{plan}(t)$；日内层以历史路径平均价值函数 $\bar{H}_t(e)$ 引导 DP 价值执行器实时调度储能，缺口按 5 倍波动电价紧急购电补齐。唯一的改动是电价口径：日前层目标函数变为

\begin{equation}
  \min_{E_{plan},\{E^{s}\}}\;
  \sum_{t=1}^{144} Price_d(t)\,E_{plan}(t)
  + \sum_{s=1}^{M}\pi_s \sum_{t=1}^{144} \kappa\, Price_d(t)\, E_{emg}^{s}(t),
  \qquad \kappa=5
  \label{eq:q4_obj}
\end{equation}

其中 $Price_d(t)$ 为规划日 $d$ 各时段的波动电价；日内层 DP 递推的即时成本同步替换为 $5\,Price_d(t)\,[r_t+\psi(x)]_{+}$，价值函数与结算口径一致更新。储能参数、容量与功率约束、跨日连续边界与终端价值引导均与问题二完全相同。

\subsection{Q4-2 求解与验证}

按问题二的流程对 334 天逐日求解：每天 0:00 生成历史情景并求解两阶段 LP 得 $E_{plan}$，构建 DP 价值函数后按真实负载、光伏逐时段执行，全部 334 天均获得最优解，总求解耗时约 $4{,}157$ 秒。对结果执行与问题二相同的验证：全部时段的储能电量保持在 $[1200, 10800]$ kWh 内、单时段充放电量不超过 $833.33$ kWh，违反次数均为 0；333 个跨日过渡点储电量衔接误差均在容差内；首日初始电量 $6000$ kWh，年末电量 $6076.41$ kWh，相对循环目标的偏差 $1.27\%$，与问题二同量级。

\subsection{Q4-2 结果分析}

\subsubsection{年度成本与结构}

波动电价下 334 天的年度总购电费用为 \textbf{\$14{,}777{,}634.02\$ 元}，其中计划购电费用 $13{,}481{,}079.23$ 元（$91.2\%$），紧急购电费用 $1{,}296{,}554.79$ 元（$8.8\%$），日均成本 $44{,}244.41$ 元。与问题二（固定电价）的 $14{,}120{,}324$ 元相比增加 $657{,}310$ 元（$4.7\%$）：波动电价下部分时段电价高于固定水平，且价格信号逐日变化给日前计划带来额外的适应成本，但成本结构保持稳定（计划/紧急占比分别为 $91.1\%$ 与 $8.9\%$、$91.2\%$ 与 $8.8\%$），表明两阶段框架在波动价格环境下仍然有效。

能源结构方面，全年计划购电量 20{,}797{,}140.61 kWh、紧急购电量 249{,}484.25 kWh（占总购电量 $1.2\%$），有效购电价格 $0.7021$ 元/kWh；紧急购电事件共 2{,}765 次（日均 8.3 次），与问题二的 246{,}579 kWh 基本相当。

\subsubsection{购电成本的时序特征}

334 天日购电成本呈现明显的季节与事件双重结构：6--7 月与 12 月日均成本最高（约 53,650、54,736 与 64,444 元/天），对应波动电价的夏季、冬季高位区间；4--5 月最低（约 32,560、31,078 元/天）。全年最高日成本出现在 6 月 1 日（119,768.83 元），最低为 4 月 25 日（9,379.79 元），日际成本波动主要跟随波动电价与光伏出力的联合变化。与问题二相比，成本的时间分布从“由光伏季节主导”转为“由电价季节与光伏季节共同主导”，这正是波动电价环境下的关键策略差异。

\subsubsection{指定日期的购电结果}

表~\ref{tab:q4_days}汇总四个指定日期（2025.3.20、2025.6.21、2025.9.23、2025.12.21）在波动电价下的购电成本与储能运行结果。夏至日（6.21）总成本仅 18,130.89 元且无紧急购电，对应光伏出力高峰与相对温和的电价水平；冬至日（12.21）总成本达 73,147.46 元，对应光伏出力全年低谷与冬季电价高位；春分与秋分日介于两者之间，且均触发了一定规模的紧急购电。

\begin{table}[H]
  \centering
  \caption{波动电价下指定日期的购电结果}
  \label{tab:q4_days}
  \small
  \begin{tabular}{lcccc}
    \toprule
    指标 & 2025.3.20 & 2025.6.21 & 2025.9.23 & 2025.12.21 \\
    \midrule
    总购电成本(元) & 42,863.04 & 18,130.89 & 45,853.21 & 73,147.46 \\
    其中：计划购电(元) & 41,399.89 & 18,130.89 & 43,201.84 & 73,147.46 \\
    其中：紧急购电(元) & 1,463.15 & 0.00 & 2,651.37 & 0.00 \\
    日初储电量(kWh) & 1,546.49 & 2,849.80 & 2,627.10 & 1,471.02 \\
    日末储电量(kWh) & 1,200.00 & 3,287.33 & 1,716.98 & 1,589.15 \\
    \bottomrule
  \end{tabular}
\end{table}

\subsection{Q4-3：波动电价下的滚动调整购电模型}

\subsubsection{模型架构}

Q4-3 对应“问题三模型 + 波动电价”：日前计划、6:00/12:00/18:00 三次预报调整与违约计费框架均继承问题三，计划购电费、调整费用与紧急购电费用的计价基准全部替换为附件4 的波动电价序列 $Price_d(t)$。在实现中，紧急购电保持与问题二一致的“补偿项”定位：优化模型求解计划与调整购电量及储能调度，紧急购电不作为优化变量，而是实时执行阶段供电不足时的被动补充（按 $5$ 倍波动电价结算）。该架构与 Q4-2 的两阶段补偿结构一致。

\subsubsection{求解与年度成本}

对 334 天滚动求解，全部天数获得最优解，总运行时长约 $62.5$ 分钟（平均每日 $11.2$ 秒）。全年总运营成本为 \textbf{18.07 百万元}，其中计划购电成本 17.66 百万元（$97.8\%$）、调整购电成本 0.18 百万元（$1.0\%$）、紧急购电成本 0.22 百万元（$1.2\%$）。与问题三的固定电价情形（16.49 百万元）相比，波动电价下总成本增加 1.58 百万元（$9.6\%$），成本增量主要来自电价波动风险与预测不确定性导致的非最优时段购电。

\begin{table}[H]
  \centering
  \caption{Q4-3 与问题三的成本对比}
  \label{tab:q43_vs_q3}
  \small
  \begin{tabular}{lcccc}
    \toprule
    情形 & 计划购电 & 调整购电 & 紧急购电 & 总成本 \\
    \midrule
    问题三（固定电价） & 14.97 百万元 & 0.31 百万元 & 1.21 百万元 & 16.49 百万元 \\
    Q4-3（波动电价） & 17.66 百万元 & 0.18 百万元 & 0.22 百万元 & 18.07 百万元 \\
    \bottomrule
  \end{tabular}
\end{table}

\subsubsection{紧急购电与储能运行}

紧急购电的日均占比为 \textbf{0.39\%}（最大单日占比 $11.90\%$），全年共 301 天发生紧急购电，总紧急购电量 76{,}517.96 kWh。作为对比，若将紧急购电作为优化变量直接纳入模型（v1 架构），回测显示其占比高达 $66.7\%$、总成本约 32 百万元；将紧急购电改为实时补偿项后（v2 架构），占比降至 $0.39\%$、总成本降至 18.07 百万元，验证了补偿式架构对模型行为的决定性影响。

购电量结构方面，全年总购电量 2{,}238.72 万 kWh，其中计划购电 2{,}217.07 万 kWh（$99.03\%$）、调整购电 14.00 万 kWh（$0.63\%$）、紧急购电 7.65 万 kWh（$0.34\%$）。储能全年总充电量 1{,}207.25 万 kWh、总放电量 1{,}086.45 万 kWh，两者之比约 $90\%$，与单侧充放电效率 $\eta=0.9$ 一致；实际平均购电电价为 $0.7970$ 元/kWh，较附件4 全年均价 $0.7662$ 元/kWh 高 $4.0\%$，价差主要来自调整购电的非最优时段执行与紧急购电的 5 倍惩罚。月度成本在 12 月最高（252.48 万元）、5 月最低（112.39 万元），冬季与夏季（平均约 187 万元/月）高于春秋季（平均约 133 万元/月）。

\subsection{问题四小结}

问题四在波动电价环境下重算了问题二与问题三。Q4-2 年度总成本 $\num{14777634.02}$ 元，较固定电价情形增加 $4.7\%$，成本结构保持稳定；Q4-3 年度总成本 18.07 百万元，较问题三增加 $9.6\%$，紧急购电占比 $0.39\%$。两个子问题的全部约束均零违反，求解成功率 $100\%$。结果表明，本文的“日前优化 + 日内执行”分层框架对电价波动具有良好的适应性：仅需替换价格输入，模型无需结构性修改即可迁移至波动电价场景，紧急购电的补偿式架构在波动价格下依然有效。

% 模型检验与稳健性讨论
\section{模型检验与稳健性讨论}

\subsection{问题一模型的检验小结}

问题一的求解结果已通过四级验证（第~\ref{subsec:verify}节）：输入数据完整合规；求解器返回 Optimal；144 个时段的电力平衡、储能状态转移、容量与功率约束逐项复核通过；全天能量守恒误差为 $0.0000$ kWh。此外，诊断检查显示无同时充放电时段、供需富余量仅 $0.0022$ kWh，表明最优策略在物理与经济意义上均合理。

\subsection{问题二模型的检验小结}

问题二的求解结果通过了跨日连续性、硬约束、能量效率与执行器对照四类验证（第~\ref{subsec:verify2}节）：333 个跨日过渡点的储电量衔接误差均小于 $10^{-6}$ kWh；全部 $334\times144$ 个时段的容量与功率约束零违反；储能全年循环效率 $81.00\%$，与理论往返效率一致；DP 执行器在固定计划对照实验中优于解析响应（$1.07\%$）、MPC（$1.80\%$）与固定计划基准（$14.31\%$，$p<0.0001$）。边界条件的三种设定对比进一步表明，跨日连续 + 年度循环的方案以几乎可忽略的成本代价（较无终端约束方案仅多 $0.016\%$）保证了年度运行的可循环性，模型在性能与可行性之间达到稳健平衡。

\subsection{模型稳健性讨论}

% TODO(C)：Q3--Q4 完成后，汇总各问的模型检验与稳健性讨论。
（本节内容待全部子问题求解完成后撰写。）

% 模型评价与推广
\section{模型评价与推广}

\subsection{模型优点}

（本节内容待全部子问题完成后由 C 撰写：线性规划模型的全局最优性、可复现性与求解效率；递进式模型体系对不确定性逐层扩展的适应性；严格的逐时段验证流程等。）

\subsection{模型局限}

（本节内容待全部子问题完成后由 C 撰写：未计入储能循环寿命与折旧成本\cite{reniers2018improving}；未考虑电网功率约束与电能质量要求；确定性假设在真实运行中的偏差等。）

\subsection{模型推广}

（本节内容待全部子问题完成后由 C 撰写：向多微网协同、需求响应、碳交易成本等场景的推广方向。）


% 2026 强制要求：AI 工具使用声明必须置于参考文献之前。
% 措辞由全队核对后定稿；支撑材料须含 AI工具使用详情.pdf。
\CumcmAIUsed{数学模型的辅助审查与验证、程序编写与调试、论文初稿撰写与语言润色}

\nocite{zia2018microgrids,wu2015demand,parra2016effect,hou2023multi,birge1997stochastic,garciatorres2021stochastic}
\bibliographystyle{gbt7714-numerical}
\bibliography{references}

\newpage
\appendix
% 附录一级标题格式(仅针对附录更改，不要改sty文件)
\titleformat{\section}
  {\centering\heiti\zihao{3}\bfseries}
  {附录~\thesection}
  {0.8em}
  {}

% 附录
\section{支撑材料文件列表}

支撑材料已单独压缩为 RAR 或 ZIP 文件，文件清单如表~\ref{tab:support}所示；提交前必须与实际压缩包逐项核对一致。

\begin{longtable}{@{}p{0.08\textwidth}p{0.42\textwidth}p{0.42\textwidth}@{}}
  \caption{支撑材料文件列表}
  \label{tab:support} \\
  \toprule
  序号 & 文件路径 & 内容说明 \\
  \midrule
  1 & \texttt{code/run\_q1.py} & 问题一主程序（调用 step1--step3，输出 result1.xlsx） \\
  2 & \texttt{code/step1\_data.py} & 问题一数据读取与输入验证 \\
  3 & \texttt{code/step2\_solve.py} & 问题一线性规划模型构建与 HiGHS 求解 \\
  4 & \texttt{code/step3\_deliverables.py} & 问题一结果提取、硬约束验证与交接文件生成 \\
  5 & \texttt{code/generate\_figures.py} & 问题一论文图表生成脚本 \\
  6 & \texttt{code/q2\_data\_loader.py} & 问题二数据加载与输入验证 \\
  7 & \texttt{code/q2\_scenario\_builder.py} & 问题二历史联合残差情景构造 \\
  8 & \texttt{code/q2\_dp\_executor.py} & 问题二动态规划价值执行器 \\
  9 & \texttt{code/q2\_full\_year\_backtest.py} & 问题二 334 天回测主程序 \\
  10 & \texttt{code/q2\_generate\_all\_results.py} & 问题二结果汇总与 result2.xlsx 生成 \\
  11 & \texttt{code/q2\_generate\_cn\_figures.py} & 问题二论文图表生成脚本 \\
  12 & \texttt{code/q4\_Q4\_2\_solver\_v2.py} & 问题四（波动电价）主求解器 \\
  13 & \texttt{code/q4\_data\_loader.py} & 附件1--4 数据加载与验证 \\
  14 & \texttt{code/q4\_scenario\_builder.py} & 波动电价下历史情景构造 \\
  15 & \texttt{code/q4\_dp\_executor.py} & 波动电价下动态规划价值执行器 \\
  16 & \texttt{code/q4\_generate\_result4\_2.py} & result4-2.xlsx 生成脚本 \\
  17 & \path{code/q4_fill_template_correct.py} & 结果表格修正与填写 \\
  18 & \texttt{code/q4\_generate\_figures.py} & 问题四论文图表生成脚本 \\
  19 & \path{code/q3_q3_model_scenario_v1.2.py} & 问题三多阶段滚动优化主模型 \\
  20 & \path{code/q3_diagnostic_334days_run.py} & 问题三 334 天诊断运行脚本 \\
  21 & \path{code/q3_curtailment_analysis.py} & 问题三弃光分析脚本 \\
  22 & \path{code/q3_final_energy_audit.py} & 问题三能量守恒审计脚本 \\
  23 & \path{code/q3_energy_balance_theory.py} & 问题三能量平衡理论核算 \\
  24 & \path{code/q3_strict_energy_diagnosis.py} & 问题三严格能量诊断 \\
  25 & \texttt{AI工具使用详情.pdf} & AI 工具使用情况说明（按 2026 年规定） \\
  \bottomrule
\end{longtable}

% TODO(C)：Q4-3 的代码随对应问题求解完成逐步补充。

\section{问题一源程序}

建模所用全部源程序完整、可运行，与论文结果一致。

\lstinputlisting[language=Python,caption={问题一主程序 run\_q1.py}]{code/run_q1.py}
\lstinputlisting[language=Python,caption={数据读取与输入验证 step1\_data.py}]{code/step1_data.py}
\lstinputlisting[language=Python,caption={LP 模型构建与求解 step2\_solve.py}]{code/step2_solve.py}
\lstinputlisting[language=Python,caption={结果提取、验证与交接 step3\_deliverables.py}]{code/step3_deliverables.py}
\lstinputlisting[language=Python,caption={论文图表生成 generate\_figures.py}]{code/generate_figures.py}

\section{问题二源程序}

\lstinputlisting[language=Python,caption={数据加载与输入验证 q2\_data\_loader.py}]{code/q2_data_loader.py}
\lstinputlisting[language=Python,caption={历史联合残差情景构造 q2\_scenario\_builder.py}]{code/q2_scenario_builder.py}
\lstinputlisting[language=Python,caption={动态规划价值执行器 q2\_dp\_executor.py}]{code/q2_dp_executor.py}
\lstinputlisting[language=Python,caption={334 天回测主程序 q2\_full\_year\_backtest.py}]{code/q2_full_year_backtest.py}
\lstinputlisting[language=Python,caption={结果汇总与输出 q2\_generate\_all\_results.py}]{code/q2_generate_all_results.py}
\lstinputlisting[language=Python,caption={图表生成 q2\_generate\_cn\_figures.py}]{code/q2_generate_cn_figures.py}

\section{问题四源程序}

\lstinputlisting[language=Python,caption={波动电价主求解器 q4\_Q4\_2\_solver\_v2.py}]{code/q4_Q4_2_solver_v2.py}
\lstinputlisting[language=Python,caption={数据加载与验证 q4\_data\_loader.py}]{code/q4_data_loader.py}
\lstinputlisting[language=Python,caption={历史情景构造 q4\_scenario\_builder.py}]{code/q4_scenario_builder.py}
\lstinputlisting[language=Python,caption={动态规划价值执行器 q4\_dp\_executor.py}]{code/q4_dp_executor.py}
\lstinputlisting[language=Python,caption={result4-2.xlsx 生成 q4\_generate\_result4\_2.py}]{code/q4_generate_result4_2.py}
\lstinputlisting[language=Python,caption={结果表格修正脚本}]{code/q4_fill_template_correct.py}
\lstinputlisting[language=Python,caption={图表生成 q4\_generate\_figures.py}]{code/q4_generate_figures.py}

\section{问题三源程序}

\lstinputlisting[language=Python,caption={多阶段滚动优化主模型 q3\_q3\_model\_scenario\_v1.2.py}]{code/q3_q3_model_scenario_v1.2.py}
\lstinputlisting[language=Python,caption={334 天诊断运行 q3\_diagnostic\_334days\_run.py}]{code/q3_diagnostic_334days_run.py}
\lstinputlisting[language=Python,caption={弃光分析 q3\_curtailment\_analysis.py}]{code/q3_curtailment_analysis.py}
\lstinputlisting[language=Python,caption={能量守恒审计 q3\_final\_energy\_audit.py}]{code/q3_final_energy_audit.py}
\lstinputlisting[language=Python,caption={能量平衡理论核算 q3\_energy\_balance\_theory.py}]{code/q3_energy_balance_theory.py}
\lstinputlisting[language=Python,caption={严格能量诊断 q3\_strict\_energy\_diagnosis.py}]{code/q3_strict_energy_diagnosis.py}


\end{document}
